\section{Conclusion}

Ce projet a permis de concevoir \AUSSI{}, un moteur générique de résolution de
problèmes basé sur la recherche dans un espace d'états. L'architecture repose
sur une séparation entre le moteur de recherche, les stratégies d'exploration,
les modes de gestion des duplications et les modèles spécifiques aux problèmes.

Le système prend en charge les labyrinthes, la navigation, le N-Puzzle et le
jeu arithmétique sans nécessiter de modification des algorithmes fondamentaux.
Les stratégies \BFS{}, \UCS{}, \Greedy{} et \AStar{} peuvent être sélectionnées
dynamiquement, tout comme les modes \TreeSearch{} et \GraphSearch{}.

Le benchmark ciblé confirme que la stratégie à privilégier dépend de
l'objectif. Sur le trajet routier testé, \UCS{} et \AStar{} trouvent le coût
minimal mesuré (881,9\,km), tandis que Greedy développe le moins de noeuds mais
renvoie un coût de 912,7\,km. Sur le 8-Puzzle moyen, \AStar{} avec Manhattan
trouve une solution de 28 coups en développant 1\,440 noeuds, contre 177\,134
avec $h=0$. Enfin, sur le labyrinthe pondéré, \GraphSearch{} réduit fortement
les développements de \UCS{} par rapport à \TreeSearch{} (15 contre 596).

Ces résultats restent limités aux instances décrites dans le protocole. Les
temps n'ont été mesurés qu'une fois par configuration, et le labyrinthe
pondéré ne distingue pas les stratégies par le coût obtenu. Des expériences
répétées sur davantage d'instances pondérées et de configurations du N-Puzzle
permettraient de renforcer les conclusions.

Les améliorations futures possibles comprennent l'ajout de nouvelles stratégies
de recherche, des heuristiques plus avancées, une interface graphique et un
module de génération automatique d'instances de problèmes.